\begin{afterword}
\summaryhead

The post begins with the brain, a lump of grey tissue that does not look dangerous. It then tells a parable. For millions of years predators and prey fought with teeth, claws, shells and venom, and no species had an ultimate weapon. Then came the ``Squishy Things,'' with no armor and weak fingers. Judged by the standards of evolution, their later feats (cutting stone, smelting metal, nuclear explosions, handling DNA, flight) would have looked absurd. So why, the post asks, could an artificial intelligence not do anything interesting over the Internet without a robot body?

People who doubt an intelligence explosion, the author says, often hear ``intelligence'' as book smarts. They say that intelligence cannot beat charisma or a gun, and ask where an AI would get money, forgetting that charisma, guns and money all come from human brains. The archetype of intelligence is a human being. We cannot see how one power could do so much, but our ignorance is a fact about us, not about intelligence, which is as real as electricity and far more powerful.

\responsehead

The post makes one point well. Judged by bodies and what bodies grow, the achievements of human beings would have looked impossible, and they came from minds. The closing point is also correct: not knowing how intelligence works is a fact about our knowledge, not evidence that intelligence is weak. The details check out, from Aristotle's view of the brain to the arithmetic of Planck times.

The trouble is in the step from the parable to its conclusion. The Squishy Things had fingers. They were weak, but every feat in the story began with hands that could pick up, shape and build. The conclusion concerns an AI that acts ``over the Internet'' without a robot body. The parable shows how far intelligence took creatures that could handle things. It does not show how far intelligence goes without a body, and the post does not say how such an AI would act on the world.

The replies to the objectors work in two ways that do not meet the objections. The first widens the word. The speakers use ``intelligence'' in the narrow sense of book smarts, which the post itself describes; the reply counts charisma, and everything else a brain does, as intelligence. Whether an AI built for the narrow kind would have the rest is exactly what the sayings question, and the wide definition answers it without argument. The second answers a question about one agent with a fact about the species. Guns and money were made by many people over long periods. The sayings ask what a single new intelligence could do in a world where guns and money already belong to others. That humanity made them does not answer this.

One sentence has been revised since the post was first published: it names the Machine Intelligence Research Institute, a name the institute took in 2013.

In short: a vivid parable about what intelligence did for a species with hands, over a long history, offered as an answer to doubts about a single machine without one, and a reply to those doubts that works by widening the word ``intelligence.''
\end{afterword}
